% !TEX root = ../../LAFROX_AnalisisFinanciero_2_20261113.tex
\section{Curva S}
\label{sec:e2-curva-s}

La caja del proyecto se mantiene positiva los 56 meses, con un mínimo de 414,40 UF en el mes 9, y cierra el contrato con 39.663,68 UF. Las cifras salen de las hojas \texttt{FLUJO CAJA NETO}, \texttt{FINANCIAMIENTO} y \texttt{CURVA S} del modelo (LafroX, 2026f), en valores netos y antes de impuestos.

La Curva S usa el flujo de caja financiado, que incluye todo el dinero que entra y sale de la cuenta del proyecto. Tiene dos curvas:
\begin{itemize}
  \item \textbf{Cobros acumulados:} el capital inicial del mes 1, los doce hitos de implementación del Formulario E-25 y las 36 mensualidades de la Operación.
  \item \textbf{Pagos acumulados:} los costos del proyecto de la sección~\ref{sec:e2-costos-precio}, los intereses del crédito y la devolución de su principal.
\end{itemize}

La distancia vertical entre ambas es el saldo de caja, el dinero disponible al cierre de cada mes. A diferencia del VAN y la TIR de la sección~\ref{sec:e2-evaluacion}, que usan el flujo sin financiamiento, este flujo muestra si la caja alcanza para pagar cada mes.

\subsection{Capital inicial y financiamiento}
\label{sec:e2-curva-financiamiento}

El capital inicial es de 18.000 UF y entra en el mes 1, con 14.400 UF de crédito bancario al 0,9\,\% mensual y 3.600 UF de los socios, en la proporción 80/20 del Formulario E-24 (Distribuidora Puelche S.A., 2026a). La Tabla~\ref{tab:e2-financiamiento} resume las condiciones.

\begin{tablalafrox}{Capital inicial y financiamiento del proyecto en UF, CLP y USD. Fuente: modelo financiero, hojas PARAMETROS y FINANCIAMIENTO.}{tab:e2-financiamiento}{Y R{2.3cm} R{3.0cm} R{2.4cm}}
  {\cab{Concepto} & \cab{UF} & \cab{CLP} & \cab{USD}}
  Crédito bancario (80\,\%) & 14.400,00 & 576.000.000 & 640.000 \\
  Aporte de los socios (20\,\%) & 3.600,00 & 144.000.000 & 160.000 \\
  \textbf{Capital inicial} & \textbf{18.000,00} & \textbf{720.000.000} & \textbf{800.000} \\
  \midrule
  Interés mensual, meses 2 a 15 & 129,60 & 5.184.000 & 5.760 \\
  Cuota mensual, meses 16 a 39 & 669,82 & 26.792.662 & 29.770 \\
  \textbf{Total de intereses} & \textbf{3.490,00} & \textbf{139.599.883} & \textbf{155.111} \\
\end{tablalafrox}

Hasta el mes 15 se pagan solo intereses, 129,60 UF al mes. Desde el mes 16, cuando se cobra H7, el crédito se paga en 24 cuotas de 669,82 UF hasta el mes 39. En total se pagan 3.490,00 UF de intereses.

El capital se fijó en lo mínimo que mantiene la caja sobre cero. Sin financiamiento, el déficit acumulado más alto es de 16.548,80 UF en el mes 9. Con los 18.000 UF se cubre ese déficit y los intereses pagados hasta entonces, y quedan 414,40 UF (\$16.575.858 o USD~18.418) de margen. Así el proyecto no usa la línea de crédito de 1.250 UF al 3\,\% mensual, que queda como respaldo si se atrasa el cobro de un hito.

\subsection{Costos e ingresos acumulados}
\label{sec:e2-curva-acumulados}

La Figura~\ref{fig:e2-curva-s} muestra los cobros y los pagos acumulados. La línea vertical marca el mes 21, cuando se cobra H12 por la aceptación final y parte la Operación.

\begin{figure}[htbp]\centering
  % !TEX root = ../../LAFROX_AnalisisFinanciero_2_20261113.tex
% Curva S del modelo financiero: cobros y pagos acumulados, en miles de UF, meses 1 a 56.
\begin{tikzpicture}
\begin{axis}[
  width=\textwidth, height=0.52\textwidth,
  font=\fontsize{9pt}{11pt}\selectfont,
  xmin=0, xmax=57, ymin=0, ymax=270,
  xtick={1,4,8,12,16,20,24,28,32,36,40,44,48,52,56},
  ytick={0,50,100,150,200,250},
  xlabel={Mes del contrato}, ylabel={Miles de UF},
  axis lines*=left, tick align=outside,
  tick label style={/pgf/number format/use comma,/pgf/number format/1000 sep={.}},
  ymajorgrids, grid style={lafroxTenue},
  legend style={at={(0.5,-0.22)},anchor=north,draw=none,legend columns=2,
    /tikz/every even column/.append style={column sep=10pt}},
  legend cell align=left,
]
\draw[lafroxGris,dashed,thin] (axis cs:20.5,0) -- (axis cs:20.5,270);
\node[anchor=north east,lafroxGris,font=\fontsize{9pt}{11pt}\selectfont] at (axis cs:20.2,262) {Implementación};
\node[anchor=north west,lafroxGris,font=\fontsize{9pt}{11pt}\selectfont] at (axis cs:20.8,262) {Operación};
\addplot[lafroxNaranja,thick,mark=none] coordinates {(1,19.000) (2,27.077) (3,27.077) (4,34.145) (5,34.145) (6,44.241) (7,44.241) (8,44.241) (9,44.241) (10,54.337) (11,54.337) (12,64.434) (13,69.482) (14,74.530) (15,74.530) (16,84.626) (17,94.723) (18,104.819) (19,109.867) (20,109.867) (21,123.514) (22,127.064) (23,130.614) (24,134.164) (25,137.714) (26,141.264) (27,144.814) (28,148.364) (29,151.914) (30,155.464) (31,159.014) (32,162.564) (33,166.114) (34,169.664) (35,173.214) (36,176.764) (37,180.314) (38,183.864) (39,187.414) (40,190.964) (41,194.514) (42,198.064) (43,201.614) (44,205.164) (45,208.714) (46,212.264) (47,215.814) (48,219.364) (49,222.914) (50,226.464) (51,230.014) (52,233.564) (53,237.114) (54,240.664) (55,244.214) (56,247.764)};
\addlegendentry{Cobros acumulados}
\addplot[lafroxFuerte,thick,mark=none] coordinates {(1,2.021) (2,17.933) (3,20.486) (4,23.734) (5,29.338) (6,33.717) (7,37.968) (8,40.742) (9,43.301) (10,45.702) (11,47.881) (12,49.398) (13,54.284) (14,59.101) (15,63.962) (16,69.860) (17,74.996) (18,79.171) (19,85.001) (20,90.896) (21,96.458) (22,101.169) (23,105.040) (24,108.762) (25,112.916) (26,116.606) (27,120.240) (28,123.922) (29,127.552) (30,131.212) (31,134.879) (32,138.907) (33,142.556) (34,146.096) (35,150.070) (36,153.636) (37,157.647) (38,161.166) (39,164.707) (40,167.548) (41,170.327) (42,173.162) (43,176.015) (44,179.216) (45,181.983) (46,184.663) (47,187.812) (48,190.546) (49,194.003) (50,196.691) (51,199.368) (52,202.091) (53,204.721) (54,207.469) (55,210.189) (56,213.372)};
\addlegendentry{Pagos acumulados}
\end{axis}
\end{tikzpicture}

  \caption{Curva S del proyecto: cobros y pagos acumulados, en miles de UF. Fuente: modelo financiero, hoja CURVA S.}
  \label{fig:e2-curva-s}
\end{figure}

Los cobros suben en escalones con cada hito y luego de forma pareja durante la Operación. Los pagos suben desde el primer mes, con un salto en el mes 2 por la compra de equipos. La Tabla~\ref{tab:e2-curva-hitos} da los valores en los meses donde cambian las curvas.

{\renewcommand{\LFXfiletefila}{\noalign{{\color{lafroxLinea}\hrule height 0.3pt}\nobreak}}%
\begin{tablalafrox}{Cobros y pagos acumulados y saldo de caja en meses seleccionados, en UF. Fuente: modelo financiero, hoja CURVA S.}{tab:e2-curva-hitos}{Y R{2.8cm} R{2.8cm} R{2.8cm}}
  {\cab{Mes} & \cab{Cobros} & \cab{Pagos} & \cab{Saldo de caja}}
  1, capital inicial & 18.000,00 & 2.020,66 & 15.979,34 \\
  2, compra del equipamiento & 26.210,41 & 17.925,38 & 8.285,03 \\
  9, saldo mínimo & 43.657,52 & 43.243,12 & 414,40 \\
  10, hito H4 & 53.920,52 & 45.637,11 & 8.283,42 \\
  16, primera cuota del crédito & 84.709,55 & 70.049,46 & 14.660,09 \\
  21, H12 e inicio de la Operación & 125.013,39 & 97.722,09 & 27.291,30 \\
  39, última cuota del crédito & 203.913,15 & 176.245,72 & 27.667,42 \\
  56, fin del contrato & 278.429,58 & 238.765,90 & 39.663,68 \\
\end{tablalafrox}}

Al mes 56 se han cobrado 278.429,58 UF, que son 260.429,58 UF del contrato y 18.000 UF de capital. Se han pagado 238.765,90 UF, entre 220.875,90 UF de costos, 3.490,00 UF de intereses y 14.400 UF de devolución del crédito. Sin el capital ni su devolución, los ingresos del contrato suman 260.429,58 UF y los costos con intereses 224.365,90 UF. Los cobros van siempre sobre los pagos, y donde más se acercan es en el mes 9.

\subsection{Análisis del flujo de caja mensual}
\label{sec:e2-curva-flujo-mensual}

En la implementación, 11 de los 20 meses tienen flujo negativo: 1, 2, 3, 5, 7, 8, 9, 11, 15, 19 y 20. El peor es el mes 2, con la compra de equipos por 12.975,08 UF, de la que H1 paga 8.210,41 UF y el capital inicial el resto. La caja llega a 414,40 UF en el mes 9 y sube a 8.283,42 UF en el mes 10 con el cobro de H4. Sin financiamiento, el acumulado es positivo desde el mes 17.

En la Operación entra una cuota de 4.383,32 UF al mes contra un costo promedio de 3.723,37 UF, y el mes 21 suma H12, de 10.263,01 UF. Los 36 meses dejan flujo positivo antes del crédito. Con la cuota del crédito, el mes queda en negativo en los meses 22, 25, 32, 35 y 37: en el 25 y el 37 se renuevan garantías y seguros, y en el 22, el 32 y el 35 hay más horas programadas. La caja acumulada sigue positiva en todos ellos, y desde el mes 40, ya pagado el crédito, el excedente se acumula.

El saldo final de 39.663,68 UF es el margen del proyecto, 39.553,68 UF, más los 3.600 UF de los socios y menos los 3.490,00 UF de intereses.
